
\documentclass[12pt]{amsart}
\pagestyle{plain} 
\usepackage[utf8]{inputenc}
\usepackage{setspace}
\usepackage[english]{babel}
\usepackage{mathabx}
\usepackage{framed}
\usepackage{soul}
\usepackage[all]{xy}
\usepackage[margin=1in]{geometry}
\usepackage{arydshln}
\usepackage[colorinlistoftodos]{todonotes}
\usepackage{pb-diagram}
\usepackage[mathscr]{euscript}
\usepackage{multicol}
\usepackage{stmaryrd} %\mapsfrom
\usepackage{empheq} %for labeling equations individually in cases
\usepackage{tikz-cd}
\tikzset{
  symbol/.style={
    draw=none,
    every to/.append style={
      edge node={node [sloped, allow upside down, auto=false]{$#1$}}}
  }
}


\usepackage{amsmath,amsfonts,amssymb,amsthm,mathrsfs,latexsym,mathtools}
\usepackage{commath}
\linespread{1}
\usepackage[T1]{fontenc}
\usepackage{hyperref}
\usepackage{caption} %figure
\usepackage{subcaption} %figure

%\usepackage{lipsum}

\DeclareMathAlphabet{\mathbbmsl}{U}{bbm}{m}{sl}

\usetikzlibrary{matrix}

\newcommand{\C}{\mathbb{C}} %newly added on Aug 30, 2021 
\newcommand{\Z}{\mathbb{Z}}
\newcommand{\R}{\mathbb{R}}
\newcommand{\D}{\mathbb D}
\newcommand{\bbQ}{\mathbb{Q}}
\newcommand{\Hom}{\textup{Hom}}
\newcommand{\Ext}{\textup{Ext}}
\newcommand{\Id}{\mathrm{Id}}
\newcommand{\Gr}{\mathrm{Gr}}
\newcommand{\disc}{\mathrm{disc}}
\newcommand{\Hv}{H_{\textup{van}}}
\newcommand{\Hp}{H_{\textup{prim}}}
\newcommand{\prim}{\textup{prim}}
\newcommand{\van}{\textup{van}}
\newcommand{\Res}{\textup{Res}}
\newcommand{\Pic}{\textup{Pic}}
\newcommand{\Div}{\textup{Div}}
\newcommand{\Aut}{\textup{Aut}}
\newcommand{\tr}{\textup{tr}}
\newcommand{\Sym}{\textup{Sym}}
\newcommand{\inv}{\textup{inv}}
\newcommand {\OO} {{\mathcal O}}
\newcommand{\PP} {\mathbb{P}}
\newcommand{\cS}{\mathcal{S}}
\renewcommand{\S}{\cS}
\newcommand{\E}{\mathcal{E}}
\newcommand{\Kum}{\textup{Kum}}



\usepackage{amscd} % Aug 19, 2021

\usepackage{enumitem} %Sep 12, 2021
\usepackage{comment}%July 10, 2025


\newtheorem{theorem}{Theorem}[section]
\newtheorem{definition}[theorem]{Definition}
\newtheorem{proposition}[theorem]{Proposition}
\newtheorem{corollary}[theorem]{Corollary}
\newtheorem{question}[theorem]{Question}
\newtheorem{example}[theorem]{Example}
\newtheorem{lemma}[theorem]{Lemma}
\newtheorem{conjecture}[theorem]{Conjecture}
\newtheorem{problem}[theorem]{Problem}
\newtheorem{remark}[theorem]{Remark}
\newtheorem{claim}[theorem]{Claim}
\newtheorem{notation}[theorem]{Notation}
\newtheorem{speakernote}[theorem]{Speaker's Note}

\newtheorem{thm}{Theorem}
\newenvironment{thmbis}[1]
  {\renewcommand{\thethm}{\ref{#1}$'$}%
   \addtocounter{thm}{-1}%
   \begin{thm}}
  {\end{thm}}

%Yilong added on 2022.07.29
\tikzset{commutative diagrams/.cd,
mysymbol/.style={start anchor=center,end anchor=center,draw=none}
}
\newcommand\MySymb[2][(-1)]{%
  \arrow[mysymbol]{#2}[description]{#1}}


\title{Comments to AI (gpt5.5)'s proof of LICT$_\Z$ for $H^1$}
\author{Yilong Zhang}
\date{May 17, 2026}

\begin{document}

\maketitle


\subsection{Background of the math problem}


The invariant cycle theorem was proved by Clemens and Schmid in the 70's, and it states that for a semistable family of complex projective varieties, the invariant class on a nearby smooth fiber extends to the total space. The theorem works over rational coefficients, and it is in general not known whether the theorem holds over integral coefficients. In \cite{AGZ_v1}, Arapura-Greer-Zhang showed that it is false for $H^2$. 

A natural question to ask is whether the integral version of the local invariant cycle theorem holds in degree one (LICT$_\Z$ for $H^1$). In a subsequent paper \cite{AGZ_v2}, they showed that this is true. The proof uses Mumford's construction for degeneration of abelian varieties. Before the proof was released on the Internet, AI had produced an independent (and different) proof \cite{AI}. It is digested below.

\subsection{Prompt}
Prove or disprove the following statement:

(LICT$_\Z$ for $H^1$) \textit{Let \( X \to B \) be a one-parameter family of complex projective varieties over a holomorphic disk \( B \), and the fiber of 0 is a simple normal crossing divisor. Let \( t \in B \setminus \{0\} \), and denote the fiber over \( t \) by \( X_t \). Let \( T \) be the monodromy operator acting on \( H^1(X_t, \mathbb{Z}) \). Then the natural restriction map
\[
H^1(X, \mathbb{Z}) \to H^1(X_t, \mathbb{Z})^T
\]
to the \( T \)-invariant submodule is surjective.}

\subsection{Difficulty of the problem} The problem is of moderate difficulty.

\subsection{AI's proof} The AI's output can be found in \cite{AI}, here we give an exposition. 

The Clemens-Schmid sequence is sliced by the Wang sequence and the exact sequence of the pair $(X,X^*)$, both are exact integrally. Take an invariant class $\alpha\in H^1(X_t,\Z)^T$. By the exactness of the Wang sequence, it lifts to a class $\tilde{\alpha}\in H^1(X^*,\Z)$. The key step is to show that

\begin{lemma}(\cite[Step 2]{AI})\label{lemma_1}
    The coboundary map $H^1(X^*,\Z)\to H^2(X,X^*,\Z)$ is identified with the residue map, which takes each $\eta\in H^1(X^*,\Z)$ to $\textup{Res}_{D_i}(\eta)$ for each component $D_i$ of the central fiber $X_0=\cup_iD_i$.
\end{lemma}
\begin{lemma}(\cite[Step 2]{AI})\label{lemma_2}
    For a given invariant class $\alpha\in H^1(X_t,\Z)^T$, the lift $\tilde{\alpha}\in H^1(X^*,\Z)$ satisfies $\textup{Res}_{D_i}(\tilde{\alpha})=\textup{Res}_{D_j}(\tilde{\alpha})$ for all $i,j$.
\end{lemma}

Then, by Lemma \ref{lemma_2}, up to adjusting by a class pulled back from $H^1(B^*,\Z)$, $\tilde{\alpha}$ lies in the kernel of the coboundary map $H^1(X^*,\Z)\to H^2(X,X^*,\Z)$. Then by the exactness of the cohomology of the pair $(X,X^*)$, it lifts to $H^1(X,\Z)$. Q.E.D.


\subsection{Comments to AI's proof}
The method in AI's proof that interprets the coboundary map $H^1(X^*)\to H^2(X,X^*)$ as the residue map (cf. Lemma \ref{lemma_1}) is original. This actually holds in higher degrees by working with Griffiths' residue, but in degree one, this interpretation is the classical residue and is particularly useful for LICT$_\Z$, since one has Lemma \ref{lemma_2}, which is essential in the proof.

The proof is elementary and elegant. (Lemma \ref{lemma_1} is essentially residue theory in complex analysis, and \ref{lemma_2} is essentially reduced to Picard--Lefschetz theory for curves.) It provides a different point of view on the integral version of the invariant cycle theorem in degree one.




\bibliographystyle{plain}
\bibliography{bibfile}


\end{document}